%============================================================
%  Chapter 9 – AI Integration
%============================================================
\chapter{AI Integration}

\section{Overview}

\projectname{} integrates the \textbf{OpenAI API} \cite{openai2024} to provide an AI-powered project scoping assistant. The assistant is accessible to clients from their dashboard under the route \texttt{/client/ai-assistant}. Its primary purpose is to help non-technical clients who know what they want but lack the technical vocabulary to specify it in a form understandable to developers.

\section{System Design}

The AI assistant operates as a \textbf{structured-output chat interface}. The client types a natural language description of their project (e.g., \textit{``I need a mobile-friendly website for my bakery that lets customers order cakes online''}), and the assistant returns a structured response that can directly populate a new project posting.

\subsection{Request Flow}

\begin{enumerate}
  \item The client sends their project description to \texttt{POST /api/ai/scope}.
  \item The backend \texttt{ai.controller.ts} constructs a system prompt and user message, then calls \texttt{openai.chat.completions.create()} with the \texttt{gpt-4o} model.
  \item The OpenAI response is parsed and the structured output is returned as JSON.
  \item An \texttt{AiScopingSession} document is created in the database storing both the raw input and structured output.
  \item If the client proceeds to post the project, the session is linked to the resulting \texttt{Project} via \texttt{projectId} and the project is flagged \texttt{isAiScoped: true} with an \texttt{aiScopeSummary} text.
\end{enumerate}

\subsection{System Prompt Design}

The system prompt instructs the model to behave as a freelance project scoping expert and to respond only with a valid JSON object conforming to a specific schema:

\begin{lstlisting}[language=TypeScript, caption={AI System Prompt Structure}]
const systemPrompt = `
You are an expert freelance project scoping assistant.
Given a client's project description, return a JSON object with:
{
  "category": string,         // e.g. "Web Development"
  "subCategory": string,      // e.g. "Frontend"
  "budgetRange": {
    "min": number,
    "max": number
  },
  "projectSize": "SMALL" | "MEDIUM" | "LARGE",
  "skills": string[],         // recommended skills
  "summary": string,          // 2-3 sentence project overview
  "estimatedTimeline": string // e.g. "2-4 weeks"
}
Only respond with the JSON object. No explanations.
`;
\end{lstlisting}

\section{AiScopingSession Model}

\begin{table}[H]
\centering
\caption{AiScopingSession Model Fields}
\begin{tabularx}{\textwidth}{llX}
\toprule
\textbf{Field} & \textbf{Type} & \textbf{Description} \\
\midrule
\texttt{id}        & ObjectId & Primary key. \\
\texttt{userId}    & ObjectId & The client who initiated the session. \\
\texttt{input}     & String   & The raw natural language description provided by the client. \\
\texttt{output}    & Json     & The structured JSON response from the OpenAI API. \\
\texttt{projectId} & ObjectId? & Linked to the resulting \texttt{Project} if the client proceeded to post. \\
\texttt{createdAt} & DateTime & Session timestamp. \\
\bottomrule
\end{tabularx}
\end{table}

\section{Frontend AI Assistant Page}

The \texttt{AIAssistant.tsx} page provides a conversational chat-like UI for submitting project descriptions. Key UI elements:

\begin{itemize}
  \item A large \textbf{text area} for the project description.
  \item A \textbf{``Scope My Project''} submit button that triggers the API call.
  \item While loading: an animated spinner and a motivational message (``Analysing your project...'').
  \item On success: a structured results panel showing the recommended category, budget range, skills, and project size as readable cards.
  \item A \textbf{``Post This Project''} button that pre-fills the \texttt{PostProject} multi-step form with the AI-generated data.
\end{itemize}

\section{Error Handling and Edge Cases}

\begin{itemize}
  \item If the OpenAI API returns a non-JSON response (e.g., a refusal or safety filter trigger), the controller catches the parse error and returns a 422 status with a user-friendly message.
  \item API rate limits are handled with exponential back-off retries (up to 3 attempts).
  \item The OpenAI API key is stored as an environment variable (\texttt{OPENAI\_API\_KEY}) and never exposed to the frontend.
\end{itemize}

\section{Use of GPT-4o Model}

The \texttt{gpt-4o} model was selected for this use case due to its superior instruction-following capabilities and lower latency compared to GPT-4 Turbo. Its ability to reliably produce valid JSON from complex natural-language inputs was a critical selection criterion.
